\chapter{Freie Konstruktionen und Präsentationen}

\section{Freie Gruppen und Wortnormalformen}

\begin{definition}[Wörter mit formalen Inversen]
\label{b40guide:def:group-words}
Zu einer Menge \(X\) sei \(X^{-1}\) eine disjunkte Kopie mit Symbolen
\(x^{-1}\) für \(x\in X\); formal ist \((x^{-1})^{-1}=x\).
Ein \emph{Wort} ist eine endliche Folge von Buchstaben aus
\(X\sqcup X^{-1}\), einschließlich des leeren Wortes.
Es heißt \emph{reduziert}, wenn keine benachbarten Buchstaben
\(xx^{-1}\) oder \(x^{-1}x\) vorkommen.
Eine \emph{Kürzung} entfernt ein solches Paar.
Das inverse Wort entsteht durch Umkehrung der Reihenfolge und
Invertieren aller Buchstaben.
\end{definition}

\begin{definition}[Freie Gruppe auf einer Menge]
\label{b40guide:def:free-group}
Eine \emph{freie Gruppe auf \(X\)} ist eine Gruppe \(F(X)\) mit einer
Abbildung \(\iota_X:X\to F(X)\), so daß zu jeder Gruppe \(G\) und jeder
Belegung \(f:X\to G\) genau ein Gruppenhomomorphismus
\(\widehat f:F(X)\to G\) mit \(\widehat f\circ\iota_X=f\) existiert.
Diese Forderung heißt ihre \emph{universelle Eigenschaft}.
\end{definition}

\begin{theorem}[Existenz und Normalform der freien Gruppe]
\label{b40guide:thm:free-group-normal-form}
Für jede Menge \(X\) existiert \(F(X)\). Ihre Elemente lassen sich
als reduzierte Wörter darstellen. Jedes Wort besitzt genau eine durch
Kürzungen erreichbare reduzierte Form; die Wahl der Kürzungsschritte
ändert das Ergebnis nicht.
Die Multiplikation besteht aus Aneinanderfügen und anschließendem
Kürzen. Das leere Wort ist die Einheit, das inverse Wort die Inverse.
Die Abbildung \(\iota_X\) ordnet jedem \(x\) sein Wort der Länge eins zu
und ist injektiv. Ihr Bild erzeugt \(F(X)\).
Zwei Wörter beschreiben genau dann dasselbe Gruppenelement,
wenn ihre reduzierten Formen übereinstimmen. \cite{milne}
\end{theorem}

\begin{theorem}[Eindeutigkeit und Auswertung]
\label{b40guide:thm:free-group-universality}
Sind \((F,\iota)\) und \((F',\iota')\) freie Gruppen auf derselben Menge
\(X\), so gibt es genau einen Isomorphismus \(u:F\to F'\) mit
\(u\circ\iota=\iota'\).
Für eine Belegung \(f:X\to G\) ist ihr universeller Homomorphismus durch
\[
 \widehat f(x_1^{\varepsilon_1}\cdots x_n^{\varepsilon_n})
 =f(x_1)^{\varepsilon_1}\cdots f(x_n)^{\varepsilon_n},
 \qquad \varepsilon_i\in\{1,-1\},
\]
gegeben; beim leeren Wort ist der Wert \(1\).
Die freie Gruppe auf der leeren Menge ist die triviale Gruppe.
\cite{milne}
\end{theorem}

\section{Gruppen durch Erzeuger und Relationen}

\begin{definition}[Präsentation und Relatoren]
\label{b40guide:def:presentation}
Für eine Menge \(X\) und \(R\subseteq F(X)\) bezeichnet
\[
 \langle X\mid R\rangle
 :=F(X)/\langle\!\langle R\rangle\!\rangle_{F(X)}
\]
die durch \(X\) und \(R\) \emph{präsentierte Gruppe}; der Nenner ist der
Normalabschluß von \(R\) in \(F(X)\).
Die Elemente von \(R\) heißen \emph{Relatoren} und werden gleich \(1\)
gesetzt. Eine Schreibweise \(u=v\) steht für den Relator \(uv^{-1}\).
Eine \emph{Präsentation einer Gruppe \(G\)} umfaßt einen Isomorphismus
\(\langle X\mid R\rangle\cong G\) mit den dadurch festgelegten
Erzeugerbildern. Die Mengen \(X\) und \(R\) dürfen beliebig groß sein.
\end{definition}

\begin{theorem}[Jede Gruppe besitzt eine Präsentation]
\label{b40guide:thm:every-group-presentation}
Für jede erzeugende Belegung \(f:X\to G\), also
\(\langle f(X)\rangle=G\), gilt
\[
 G\cong F(X)/\ker\widehat f.
\]
Insbesondere ist jede Gruppe ein Quotient einer freien Gruppe.
Eine kanonische Präsentation verwendet ein eigenes Symbol \(x_g\)
für jedes \(g\in G\):
\[
 G\cong
 \bigl\langle x_g\ (g\in G)\ \bigm|\
 x_1=1,\quad x_gx_h=x_{gh}\ (g,h\in G)\bigr\rangle.
\]
Der Isomorphismus sendet \(x_g\) auf \(g\). \cite{milne}
\end{theorem}

\begin{theorem}[Abbildungsprinzip für Präsentationen]
\label{b40guide:thm:presentation-mapping}
Sei \(P=\langle X\mid R\rangle\), und sei \(\pi:F(X)\to P\)
die Quotientenabbildung. Eine Belegung \(f:X\to G\) setzt sich genau
dann zu einem Homomorphismus \(\varphi:P\to G\) fort, wenn
\[
 \widehat f(r)=1\qquad(r\in R).
\]
In diesem Fall ist die Fortsetzung eindeutig und erfüllt
\(\varphi\circ\pi=\widehat f\).
Sie ist surjektiv genau dann, wenn \(f(X)\) die Gruppe \(G\) erzeugt,
und injektiv genau dann, wenn
\(\ker\widehat f=\langle\!\langle R\rangle\!\rangle_{F(X)}\).
\cite{milne}
\end{theorem}

\begin{theorem}[Abgeleitete Relationen und benannte Erzeuger]
\label{b40guide:thm:presentation-changes}
Seien \(R\subseteq F(X)\) und
\(S\subseteq\langle\!\langle R\rangle\!\rangle_{F(X)}\).
Dann induziert die Identität auf \(X\) einen Isomorphismus
\[
 \langle X\mid R\rangle\cong\langle X\mid R\cup S\rangle.
\]
Sei ferner \(Y\) eine zu \(X\) disjunkte Menge neuer Symbole,
und sei jedem \(y\in Y\) ein Wort \(w_y\in F(X)\) zugeordnet.
Es gilt
\[
 \langle X\mid R\rangle
 \cong
 \bigl\langle X\sqcup Y\mid R,\ y=w_y\ (y\in Y)\bigr\rangle.
\]
Unter dieser Isomorphie entsprechen sich die ursprünglichen Erzeuger;
jedes neue Symbol entspricht dem durch \(w_y\) bezeichneten Element.
Die umgekehrten Operationen entfernen solche Relationen bzw. solche
zusätzlich eingeführten Erzeuger. \cite{milne}
\end{theorem}

\section{Freie Produkte beliebiger Gruppen}

\begin{definition}[Freies Produkt]
\label{b40guide:def:free-product}
Das \emph{freie Produkt} einer mengenindizierten Familie
\((G_i)_{i\in I}\) ist eine Gruppe \(P=\mathop{\ast}_{i\in I}G_i\)
mit Homomorphismen \(j_i:G_i\to P\), für die gilt:
Zu jeder Gruppe \(K\) und jeder Familie von Homomorphismen
\(f_i:G_i\to K\) gibt es genau einen Homomorphismus \(f:P\to K\)
mit \(f\circ j_i=f_i\) für alle \(i\in I\).
Für zwei Faktoren schreibt man \(G\ast H\). \cite{mcmullen}
\end{definition}

\begin{theorem}[Existenz, Eindeutigkeit und Normalform freier Produkte]
\label{b40guide:thm:free-product-normal-form}
Jede Familie \((G_i)_{i\in I}\) besitzt ein freies Produkt, eindeutig
bis auf genau einen mit allen \(j_i\) verträglichen Isomorphismus.
Die \(j_i\) sind injektiv. Jedes Element besitzt genau eine Form
\[
 j_{i_1}(g_1)\cdots j_{i_n}(g_n),\qquad
 g_k\in G_{i_k}\setminus\{1\},\quad i_k\ne i_{k+1}.
\]
Die Einheit entspricht \(n=0\). Multipliziert wird durch Aneinanderfügen,
Zusammenfassen benachbarter Faktoren aus derselben Gruppe und Entfernen
von Einheiten. Die Inverse entsteht durch Umkehrung und Invertieren.
Das freie Produkt der leeren Familie ist trivial. \cite{mcmullen}
\end{theorem}

\begin{theorem}[Präsentation eines freien Produkts]
\label{b40guide:thm:free-product-presentation}
Haben die Gruppen \(G_i\) Präsentationen \(\langle X_i\mid R_i\rangle\)
mit paarweise disjunkten Erzeugermengen, dann gilt
\[
 \mathop{\ast}_{i\in I}G_i
 \cong
 \left\langle\ \bigcup_{i\in I}X_i\ \middle|\
                    \bigcup_{i\in I}R_i\ \right\rangle.
\]
Hierbei werden die Relatoren aus \(F(X_i)\) in
\(F(\bigcup_iX_i)\) gelesen. Die Isomorphie erhält die vorgegebenen
Erzeugerbilder. \cite{milne,mcmullen}
\end{theorem}

\begin{theorem}[Vom freien zum direkten Produkt]
\label{b40guide:thm:free-to-direct-product}
Die Vorschriften \(j_G(g)\mapsto(g,1)\) und
\(j_H(h)\mapsto(1,h)\) bestimmen einen surjektiven Homomorphismus
\(q:G\ast H\to G\times H\). Sein Kern ist
\[
 \ker q=
 \left\langle\!\left\langle
 j_G(g)j_H(h)j_G(g)^{-1}j_H(h)^{-1}
 \;\middle|\; g\in G,\ h\in H
 \right\rangle\!\right\rangle_{G\ast H}.
\]
Damit entsteht \(G\times H\) aus \(G\ast H\) genau durch die zusätzlichen
Relationen, daß die beiden eingebetteten Faktoren elementweise
kommutieren. \cite{milne,mcmullen}
\end{theorem}
